\documentclass[12pt]{article}
\usepackage{amsmath}
\usepackage[margin=1in]{geometry}
\usepackage{fontspec}
\usepackage{polyglossia}
\setdefaultlanguage{thai}
\setmainfont{Noto Sans Thai}
\title{คณิตศาสตร์เบื้องหลังการเคลื่อนที่ในเกม 2D}
\author{Prompt Lab}
\date{2026}
\begin{document}
\maketitle
\begin{abstract}
บทความสั้นนี้อธิบายว่าตำแหน่ง ความเร็ว และช่วงเวลาเชื่อมโยงกันอย่างไรใน Game loop และเหตุใดการตรวจการชนจึงสามารถอธิบายด้วยกรอบสี่เหลี่ยมและอสมการได้
\end{abstract}
\section{การเคลื่อนที่ตามเวลา}
เมื่อผู้เล่นมีตำแหน่งเริ่มต้น $(x_t,y_t)$ และความเร็ว $(v_x,v_y)$ ตำแหน่งในเฟรมถัดไปคำนวณได้จาก
\[
x_{t+1}=x_t+v_x\Delta t,\qquad y_{t+1}=y_t+v_y\Delta t.
\]
\section{เงื่อนไขการชน}
วัตถุสองชิ้นชนกันเมื่อกรอบสี่เหลี่ยมของทั้งสองทับซ้อนกันในแกน $x$ และแกน $y$ ตัวอย่างเงื่อนไขในแกน $x$ คือ
\[
x_1<x_2+w_2\;\land\;x_1+w_1>x_2.
\]
ในโค้ด JavaScript สามารถนำแนวคิดนี้ไปใช้ในฟังก์ชัน \texttt{overlaps()} และตรวจแกน $y$ เพิ่มเติมก่อนตัดสินว่าผู้เล่นชนเหรียญหรือสิ่งกีดขวาง
\end{document}
